\documentclass[aspectratio=169]{beamer}
\usepackage{kotex}

% 목차의 section과 subsection을 같은 서식으로 표시
\setbeamerfont{section in toc}{family=\sffamily,size=\normalsize,series=\mdseries,shape=\upshape}
\setbeamerfont{subsection in toc}{family=\sffamily,size=\small,series=\mdseries,shape=\upshape}
\setbeamercolor{section in toc}{fg=structure.fg}
\setbeamercolor{subsection in toc}{fg=structure.fg}
\setbeamertemplate{section in toc}{%
  \leavevmode
  \leftskip=1.8em
  \llap{\raisebox{0.1ex}{\small\textbullet}\hskip0.7em}%
  \inserttocsection\par
  \vspace{0.25em}%
}

\title{Chomsky(1959), A Review of B. F. Skinner's Verbal Behavior}
\subtitle{26-2 언어습득}
\author{언어학과 이승헌}
\date{2026.09.14}

\begin{document}

% 표지
\begin{frame}
  \titlepage
\end{frame}

% 목차
\begin{frame}{Table of Contents}
  \tableofcontents
\end{frame}

\section{자극 \textit{stimulus}}
\subsection{Skinner의 자극 개념}

\begin{frame}[t]{Skinner의 자극(\textit{stimulus}) 개념}
    \vspace{1cm}
  \begin{columns}[T,onlytextwidth]
    \begin{column}[t]{0.48\textwidth}
      \begin{minipage}[t][3cm][t]{\linewidth}
        \textbf{자극}

        \medskip
        어떤 환경과 어떤 행동이 서로 법칙적으로 연관될 때,
        \textbf{‘자극과 반응’}이라는 관계에 놓인다.
      \end{minipage}
    \end{column}

    \begin{column}[t]{0.48\textwidth}
      \begin{minipage}[t][3cm][t]{\linewidth}
        \textbf{자극-반응 예시}

        \medskip
        음악을 듣고 “Mozart”라고 반응하는 것, 그림을 보고 “Dutch”라고 반응하는 것, etc. 
      \end{minipage}
    \end{column}
  \end{columns}

  \vspace{0.6em}

  \begin{columns}[T,onlytextwidth]
    \begin{column}{\textwidth}
      \textbf{법칙적인 관계}
      \medskip
      \\ 모든 반응은 그 물리적 대상(그림)의 또 다른 자극에 의해 통제된다고 설명. 만약, 빨간 의자를 보고 “빨강”이라고 말한다면 그 반응은 `\textbf{빨강성}'이라는 자극의 통제를 받는 것이고, “의자”라고 말한다면 대상 자체인 `\textbf{의자성}'의 통제를 받는다는 것이다.   
    \end{column}
  \end{columns}
\end{frame}

\subsection{Chomsky의 비판 - (1)}

\begin{frame}{Chomsky의 비판 - (1)}
    \begin{alertblock}{\textbf{단순한 만큼 공허한 설명이다!}}
    \medskip 속성은 필요할 때마다 얼마든지 만들어 낼 수 있기 때문.
    \end{alertblock}
    \bigskip
    \begin{itemize}
        \item 자극 속성을 찾아내기만 하면 거의 모든 반응을 설명할 수 있음.
        \item 화자가 처한 환경 속 자극을 바탕으로 화자의 언어행동을 예측할 수 없음.
        \item 개인이 물리적 대상의 어느 속성에 반응할지 통제할 수도 없음.
    \end{itemize}
\end{frame}

\section{강화 \textit{reinforcement}}
\subsection{Skinner의 강화 개념}

\begin{frame}[t]{Skinner의 강화(\textit{reinforcement}) 개념}
    \vspace{1cm}
  \begin{columns}[T,onlytextwidth]
    \begin{column}[t]{0.48\textwidth}
      \begin{minipage}[t][3cm][t]{\linewidth}
        \textbf{강화}
        
        \medskip
        \textbf{강화} 자극이 주어지면 반응의 \textbf{강도(빈도)}가 높아진다. 
      \end{minipage}
    \end{column}

    \begin{column}[t]{0.48\textwidth}
      \begin{minipage}[t][3cm][t]{\linewidth}
        \textbf{강화 예시}
        
        \medskip
        막대를 누르면 먹이가 나온다는 사실을 깨달은 쥐가 막대를 더 자주 누르게 되는 것, \textbf{청자(부모)의 반응이 아동의 언어 반응을 강화하는 것}, etc. 
      \end{minipage}
    \end{column}
  \end{columns}

  \vspace{0.6em}

  \begin{columns}[T,onlytextwidth]
    \begin{column}{\textwidth}
      \textbf{아동의 언어 행위를 강화하는 요인}
      
      \medskip
      \begin{itemize}
        \item 누군가가 말을 하지 않고 조용히 들어 주는 것
        \item 미래에 독자나 청자가 보일 것으로 기대되는 반응
        \item 실제로 아직 일어나지 않은, 상상하거나 희망하는 결과, etc.  
      \end{itemize} 
    \end{column}
  \end{columns}
\end{frame}

\subsection{Chomsky의 비판 - (2)}

\begin{frame}{Chomsky의 비판 - (2)}
  \begin{alertblock}{\textbf{강화 자극의 개념이 지나치게 넓어진다!}}
    \medskip
     어떤 행동을 지속하게 한 모든 가능한 요인—만족감, 기대, 상상, 미래의 반응—을 강화라고 부를 수 있다면, \underline{강화 이론은 어떠한 행동도 설명할 수 있다.}
    \newline
    \\ → 강화 이론을 통해 어떤 행동이 일어나지 않을지 예측하거나, “이 경우에는 강화 이론이 틀렸다”고 반증하는 것 또한 불가능하다는 것. 따라서 Chomsky는 강화 개념이 실험실 밖에서는 설명력을 잃는다고 봄(\textit{내용이 비게 된다}).
  \end{alertblock}
\end{frame}

\begin{frame}
  \centering
  \huge 감사합니다.
\end{frame}

\end{document}
